\labsection{3}{Uninformed and informed search implementation}{sec:lab03-search-code}

\begin{labproject}{One Python search toolkit}
\textbf{Coverage.} The ten examples in Lab 3: five uninformed-search patterns and five informed-search patterns.\\
\textbf{Goal.} Turn the repeated code into a small, testable search toolkit.

\begin{labmode}
The source handout already provides reference functions for BFS grid search, DFS maze search, UCS, BFS tree traversal, DFS preorder, GBFS grid search, A* grid search, A* graph search, and a greedy word ladder. The weighted-maze example reuses the A* grid pattern.
\end{labmode}

\medskip\noindent\textbf{Part A --- Uninformed functions.}\par
Place the following in one module: \texttt{bfs\_grid}, \texttt{dfs\_maze}, \texttt{ucs}, \texttt{bfs\_tree}, and \texttt{dfs\_preorder}.

\medskip\noindent\textbf{Part B --- Informed functions.}\par
Add \texttt{greedy\_bfs}, \texttt{astar\_grid}, \texttt{astar\_graph}, and \texttt{greedy\_word\_ladder}. Use the weighted-grid example to confirm that the A* grid implementation can accumulate cell costs.

\medskip\noindent\textbf{Part C --- Cross-algorithm trace.}\par
Choose one graph and instrument each compatible algorithm to print the frontier priority, popped state, path, and visited-set size.

\begin{tryit}
Create a summary table with columns: algorithm, data structure, priority expression, returned path, and returned cost (where defined). Use actual outputs from your code rather than a theoretical table copied from the notes.
\end{tryit}
\end{labproject}

\begin{labdeliverable}
Submit the search module plus a driver notebook/script that runs every Lab 3 example and prints a compact comparison table.
\end{labdeliverable}
